\section{One radix for every binomial graph}
\label{sec:common-radix}

The direct certificate contains sixteen binomial relations. Thirteen
certify the ranges of the CRT-encoded local coordinates; the other three
occur inside the two factorial graphs and the modulus-product graph.
The earlier implementation gives every binomial relation its own radix.
That repetition is unnecessary.

\begin{theorem}[Common radix]
Let $n_1,\ldots,n_s,k_1,\ldots,k_s$ be nonnegative integers and set
\[
 S=\sum_{i=1}^s n_i,\qquad Q=2^S+1.
\]
For each $i$ define $L_i=Q^{k_i}$ and $P_i=(Q+1)^{n_i}$.
Then $c_i=\binom{n_i}{k_i}$ for every $i$ if and only if there exist
nonnegative integers $h_i,r_i,u_i,v_i$ satisfying
\begin{align*}
 P_i&=L_i(Qh_i+c_i)+r_i,\\
 L_i&=r_i+u_i+1,\\
 Q&=c_i+v_i+1
 \qquad(1\le i\le s).
\end{align*}
We use the convention $\binom nk=0$ when $k>n$.
\end{theorem}

\begin{proof}
The binomial theorem gives
\[
 (Q+1)^{n_i}=\sum_{j=0}^{n_i}\binom{n_i}{j}Q^j.
\]
Every coefficient lies between zero and $2^{n_i}$, and
$2^{n_i}\le 2^S<Q$. Thus this is a carry-free base-$Q$ expansion.
The three displayed equations state that $r_i$ is the remainder on
division by $Q^{k_i}$ and that $c_i$ is the remainder of the quotient
on division by $Q$. Uniqueness of Euclidean division identifies $c_i$
with the indicated base digit. Conversely the actual digit and the
two quotient/remainder decompositions supply all four nonnegative
witnesses.
\end{proof}

The sixteen actual lower arguments have only four distinct values:
the CRT factorial argument $H$, the certificate length $t$, and the
two range cutoffs $m+1$ and $V+1$. Consequently the sixteen binomial
graphs need only
\[
 1+16+4=21
\]
power graphs: one for $2^S$, sixteen numerator powers, and four shared
denominator powers. The earlier compilation uses forty-eight powers
for those same binomial relations. The seven power calls outside them
are unchanged, giving a reduction from fifty-five to twenty-eight.

\paragraph{Why deferred arithmetic graphs are legitimate.}
The common radix may depend on an upper argument that itself uses a
factorial output. The factorial graph also uses that common radix.
This introduces a dependency cycle between graph predicates, but it
does not introduce a circular soundness argument. In any satisfying
assignment the common-radix theorem applies to all sixteen assigned
upper arguments, so every binomial output is its exact mathematical
value. All the original higher-level identities then follow.
For completeness, first choose the true outputs of the original
abstract power, binomial, and factorial certificate. These determine
all sixteen upper arguments. Choose the common radix from those
values, and only then fill its digit and power witnesses. The
gate-level arithmetic circuit remains triangular: graph outputs are
existential coordinates, not recursively evaluated program calls.

\section{One offset for every signed congruence}
\label{sec:one-offset}

The natural-number certificate implements a signed congruence by
\[
 L_i+M_iq_i^+=R_i+M_iq_i^-.
\]
A separate negative part for every quotient is unnecessary.

\begin{theorem}[Common quotient offset]
For any finite family of integer triples $(L_i,R_i,M_i)$, existence
of nonnegative $q_i^+,q_i^-$ satisfying the preceding equations is
equivalent to existence of nonnegative $Q,q_1,\ldots,q_s$ satisfying
\[
 L_i+M_iq_i=R_i+M_iQ\qquad(1\le i\le s).
\]
No sign or nonzero hypothesis on $M_i$ is required for this statement.
\end{theorem}

\begin{proof}
Given the separate pairs, choose a natural $Q$ with
$Q\ge\max_i q_i^-$. Put $q_i=q_i^++Q-q_i^-$. All $q_i$ are natural,
and substitution gives the common-offset equations. Conversely set
$q_i^+=q_i$ and $q_i^-=Q$. Both constructions are finite elementary
arithmetic operations.
\end{proof}

This transformation is applied across the entire finite certificate,
including congruences inside the power graphs and the outer CRT
evaluation congruence. It changes neither the externally visible
power/binomial/factorial outputs nor the numerical proof predicate.

\section{The positivity dependency in the shared frame}
\label{sec:positivity-order}

All twenty-eight power exponents are positive before any power theorem
is used. This follows from the literal nonnegative arithmetic graph:
$t=t_0+1$, $m=m_0+1$, the encoded CRT representatives are shifted by
one, and the factorial and common-radix exponents contain these positive
terms. Hence every shared Pell ordinate satisfies $Y_j\ge k_j>0$,
and their product is positive.

Twenty-seven power bases are also positive by those same syntactic
facts. The remaining base is the proposed factorial value $T=H!$;
its positivity must not be assumed. The corresponding factorial
certificate contains
\[
 P+c=cf+r+1,\qquad c=r+s+1.
\]
Its numerator $P=(\rho+1)^H$ is positive by one of the already
applicable power graphs. Therefore $c>0$, and $f=0$ would imply
$P+c=r+1\le c$, a contradiction. Thus $f\ge1$, proving positivity
of the final power base. Only after this two-stage argument are all
twenty-eight power outputs identified with their mathematical values.

The independent audit records the actual call indices in each stage;
it checks all exponents first and identifies the factorial numerator
and both ceiling rows by exact polynomial comparison. This avoids
using correctness of a factorial graph to justify a power theorem
whose own positivity assumption is needed to establish that factorial
graph.

\section{Literal checkpoints and verification scope}

Before the subsequent monic quadratic eliminations, the materialized
checkpoints are:
\begin{center}
\begin{tabular}{lrrrr}
\hline
Construction & Variables & Quadratic rows & Monomials & Degree\\
\hline
Inherited shared certificate & 3614 & 2954 & 15996 & 4\\
Common radix & 2108 & 1745 & 10629 & 4\\
Common radix and quotient offset & 1968 & 1717 & 10545 & 4\\
Full shared Pell frame & 1153 & 1021 & 9234 & 4\\
\hline
\end{tabular}
\end{center}
All variables in this table are nonnegative integers. Every polynomial
is supplied both as the sum of the listed quadratic residual squares
and as a complete expanded coefficient dictionary.

The builder-free audit reconstructs the natural affine restoration,
checks the new graph identities over the free integer polynomial ring,
verifies every displayed census, and independently expands the quartic.
The common-radix arithmetic controls include upper and lower boundary
cases and an undersized-radix carry counterexample. The actual shared
frame is also instantiated for the pair $2^1,3^1$, with all eighty
literal residuals checked exactly. A separate $A=2$ fixture with
indices $1,2$ checks the pooled modulus $64$ and period $32$.
These are regression checks; the unbounded theorems and the inherited
calculus-to-refutation calibration remain mathematical proof obligations
and are not replaced by finite testing.
